\documentclass{cssheet}


%--------------------------------------------------------------------------------------------------------------
% Basic meta data
%--------------------------------------------------------------------------------------------------------------

\title{Isomorphie}
\author{Prof. Dr. Christian Spannagel}
\date{\today}
\hypersetup{%
	pdfauthor={\theauthor},%
	pdftitle={\thetitle},% 
	pdfsubject={Aufgabenblatt Algebra},%
	pdfkeywords={algebra}
}

%--------------------------------------------------------------------------------------------------------------
% document
%--------------------------------------------------------------------------------------------------------------

\begin{document}
	\printtitle

\vspace*{10mm}


\textbf{Aufgabe 1 (Isomorphien):} 

\begin{enumerate}[a)]
\item Findet einen Isomorphismus von $(\mathbb{Z};+)$ auf $(2\mathbb{Z};+)$ und beweist, dass es sich um einen Isomorphismus handelt.
\item Beweist: Die Abbildung $\varphi(x)=2^x$ ist ein Isomorphismus von $(\mathbb{R};+)$ auf $(\mathbb{R}^{+};\cdot)$.
\end{enumerate}

\textbf{Aufgabe 2 (Gruppenübertragung):} 

\begin{enumerate}[a)]
\item Beweist: Ist $\varphi$ ein Isomorphimus von der Gruppe $(G_1;\cdot)$ auf die Gruppe $(G_2;*)$, so ist das Bild des Neutralelements von $(G_1;\cdot)$ das Neutralelement von $(G_2;*)$, und das Bild eines zu $a\in G$ inversen Elements ist das Inverse des Bildes von $a$ in $(G_2;*)$.
\item Beweist: Ist $(G;\cdot)$ eine Gruppe und $(M;*)$ eine beliebige Struktur, die zu $(G;\cdot)$ isomorph ist, dann ist $(M;*)$ eine Gruppe.
\end{enumerate}





\vspace{5cm}
\vspace*{3cm}

\printlicense

\printsocials



\end{document}